\chapter{Introduction}
\label{chap:introduction}
\section{Motivation}
Hyperbolic conservation laws describe the transport of quantities such as mass,
momentum, and energy. Their solutions can contain both smooth structures and
sharp fronts. Even smooth initial data can develop a discontinuity, so a
numerical method must accommodate qualitatively different behaviour within the
same calculation.

Discontinuous Galerkin (DG) methods combine finite element approximation within
each mesh element with numerical fluxes between elements. Independent local
polynomial representations allow high-order approximation without imposing
continuity across interfaces~\cite{HesthavenWarburton2008}. Near an unresolved
front, however, a polynomial approximation can develop overshoots and
undershoots. Damping these oscillations indiscriminately can also remove
smooth features that a high-order method is intended to resolve.

Projection-based damping offers a way to separate these concerns. The cell
mean can be retained while nonconstant parts of the approximation are
attenuated. The scalar OFDG formulation of Lu, Liu, and Shu~\cite{LuEtAl2021}
and its extension to hyperbolic systems~\cite{LiuEtAl2022} provide the starting
point. The OEDG formulation of Peng, Sun, and Wu~\cite{PengEtAl2024} further
motivates normalized sensing and exact integration of a frozen damping step.
A separate troubled-cell indicator can restrict the operation to selected
cells. The central issue is then how to define the damping, its sensor, and
its geometric realization consistently.

\section{Aim and research questions}
This thesis concerns a basis-independent realization of an adapted OFDG method
in the finite element library MFEM. Here, basis independence means that a
change in coefficient representation does not change the filtered finite
element function. The formulation combines physical projections,
face-averaged derivative jumps, and optional KXRCF gating. It is an adaptation
of the cited methods rather than a literal reproduction of a single paper.

The work is guided by three questions:
\begin{enumerate}
\item How can projection damping be expressed in the host finite element basis
while preserving component-wise cell averages?
\item How should projection, derivative sensing, and geometric normalization
be constructed when element maps are curved?
\item How do damping and selective activation affect accuracy and resolution
of smooth and discontinuous solutions?
\end{enumerate}
The first two questions determine the mathematical construction and its
implementation. The third requires an agreed set of numerical experiments;
its answer is deliberately left open in this draft.

\section{Scope and approach}
The construction uses a fixed mesh and a uniform solution order. It includes
scalar conservation laws and conservative systems, with compressible Euler
providing the system formulation. Static polynomial curved element maps are
considered; moving meshes and variable solution order are outside the scope.
On curved elements, differentiation generally leaves the mapped approximation
space. The method therefore projects first physical derivatives back into that
space and applies these projected operators recursively. The distinction
between this construction and exact higher differentiation is central to the
curvilinear extension.

The thesis separates three kinds of statements: identities derived from
projection algebra, implementation choices defining the actual method, and
empirical conclusions requiring numerical evidence. In particular, cell-mean
preservation by the filter does not by itself establish positivity, full-scheme
stability, or accuracy near a shock. Published analyses are used under their
stated assumptions, not transferred automatically to the adapted method.

\Cref{chap:background} introduces DG for a reader familiar with finite elements.
\Cref{chap:methods} develops the stabilization methods, curvilinear extension,
and implementation assumptions; \cref{sec:ofdg-implementation} gives
coefficient-space details. Results, Discussion, and Conclusion remain explicitly
incomplete pending agreement on the numerical experiments. Exploratory results
are retained in the separate preliminary brief.
